\begin{frame}{Identification relies on the orthogonality of idiosyncratic shocks}
    \vspace{-0.5em}
    \begin{block}{The Central Identification Assumption}
        Individual shocks $u_{it}$ are truly idiosyncratic—uncorrelated with aggregate shocks ($\eta_t, \epsilon_t$) and controls ($C_t^y, C_t^p$):
        \begin{equation}
            u_t \perp \eta_t, \epsilon_t, C_t^y, C_t^p
        \end{equation}
    \end{block}

    \vfill

    \begin{columns}[T]
        \begin{column}{0.48\textwidth}
            \textbf{Baseline Simplifications}
            \begin{description}
                \item[A1] Factor loadings $\lambda_i$ are known (relaxed in §IV.B).
                \item[A2] Controls have factor structure: $C_{it}^y = \lambda_i c_t^y$.
            \end{description}
        \end{column}
        \begin{column}{0.48\textwidth}
            \textbf{Error Structure}
            \begin{description}
                \item[A3] Homoskedastic and uncorrelated across $i$: $\mathbb{E}[u_t u_t'] = \sigma_u^2 I$.
                \item[A4] General heteroskedasticity $V^u$ (relaxed for optimality).
            \end{description}
        \end{column}
    \end{columns}

    \vfill
    \footnotesize
    \textit{Note: All random variables are assumed to have finite fourth moments and are i.i.d. across periods.}
\end{frame}
